\chapter{Về Tương Lai}
\markboth{Về Tương Lai}{About the Future}

\section{Nghĩ Về Tương Lai}
\begin{truyen}{Nghĩ Về Tương Lai}{Think About the Future}
\chuhoa{N}{ghĩ về tương lai là điều tự nhiên của tuổi học trò.}
\textit{(Thinking about the future is natural for students.)}

Em sẽ học gì, làm gì, trở thành ai, đi theo hướng nào.
\textit{(What will you study, do, become, or pursue?) }

Nghĩ về tương lai giúp em có hướng hơn cho hiện tại.
\textit{(Thinking about the future gives the present more direction.)}

Chỉ cần đừng để nó biến thành nỗi lo nuốt hết ngày hôm nay.
\textit{(Just do not let it turn into a worry that swallows today.)}
\end{truyen}

\begin{baihoc}
Tương lai đáng để nghĩ tới, nhưng không nên được nghĩ theo cách làm em tê liệt.
\textit{(The future is worth thinking about, but not in a way that paralyzes you.)}
\end{baihoc}

\begin{ghinhoanh}
\textbf{Short English to remember:}

The future should guide, not paralyze.

\end{ghinhoanh}

\begin{tuhocnhanh}
1. Vì sao nên nghĩ về tương lai? \textit{Why should you think about the future?}

2. Điều gì không nên xảy ra khi nghĩ về tương lai? \textit{What should not happen when thinking about the future?}

3. Em đang nghĩ gì về tương lai của mình? \textit{What are you thinking about your future?}
\end{tuhocnhanh}
\ngancach

\section{Tương Lai Được Xây Từ Hiện Tại}
\begin{truyen}{Tương Lai Được Xây Từ Hiện Tại}{The Future Is Built from the Present}
\chuhoa{N}{gày mai không rơi xuống từ trên trời.}
\textit{(Tomorrow does not drop from the sky.)}

Nó được ghép từ những việc em đang làm hôm nay: cách học, cách nghĩ, cách giữ thói quen, cách đối xử với người khác.
\textit{(It is built from what you do today: how you study, think, build habits, and treat others.)}

Hiểu điều này giúp học sinh bớt mơ hồ và bớt chỉ mơ mà không làm.
\textit{(Understanding this helps students become less vague and less likely to dream without action.)}

Tương lai thật luôn có gốc trong hiện tại thật.
\textit{(A real future always has roots in a real present.)}
\end{truyen}

\begin{baihoc}
Muốn tương lai đổi khác, em phải bắt đầu bằng việc thay đổi một phần hiện tại.
\textit{(If you want the future to change, begin by changing part of the present.)}
\end{baihoc}

\begin{ghinhoanh}
\textbf{Short English to remember:}

Today's choices build tomorrow.

\end{ghinhoanh}

\begin{tuhocnhanh}
1. Vì sao nói tương lai được xây từ hiện tại? \textit{Why do we say the future is built from the present?}

2. Điều này giúp học sinh bớt điều gì? \textit{What does this understanding reduce in students?}

3. Em muốn thay đổi gì trong hiện tại để giúp tương lai? \textit{What do you want to change now to help your future?}
\end{tuhocnhanh}
\ngancach

\section{Không Cần Biết Hết Ngay}
\begin{truyen}{Không Cần Biết Hết Ngay}{You Do Not Need to Know Everything Now}
\chuhoa{K}{hông phải học sinh nào cũng phải biết rất sớm mình sẽ trở thành gì.}
\textit{(Not every student must know very early exactly who they will become.)}

Có người rõ hướng sớm, có người cần thêm thời gian để khám phá.
\textit{(Some people find direction early; others need more time to discover it.)}

Điều quan trọng là vẫn học tử tế, sống tử tế và tiếp tục tìm hiểu mình.
\textit{(What matters is to keep learning well, living well, and understanding yourself.)}

Không biết hết ngay không có nghĩa là em đang đi sai.
\textit{(Not knowing everything now does not mean you are going wrong.)}
\end{truyen}

\begin{baihoc}
Tương lai không phải bài trắc nghiệm bắt buộc phải có đáp án thật sớm.
\textit{(The future is not a multiple-choice test that requires an early final answer.)}
\end{baihoc}

\begin{ghinhoanh}
\textbf{Short English to remember:}

Uncertainty does not always mean wrong direction.

\end{ghinhoanh}

\begin{tuhocnhanh}
1. Vì sao không cần biết hết ngay về tương lai? \textit{Why do you not need all the answers right now?}

2. Điều gì vẫn quan trọng dù chưa rõ hướng? \textit{What still matters even if the direction is unclear?}

3. Em có đang tự ép mình phải biết hết quá sớm không? \textit{Are you forcing yourself to know everything too early?}
\end{tuhocnhanh}
\ngancach

\section{Chuẩn Bị Từ Sớm}
\begin{truyen}{Chuẩn Bị Từ Sớm}{Prepare Early}
\chuhoa{C}{ó những điều chưa cần quyết định ngay, nhưng rất nên chuẩn bị từ sớm.}
\textit{(Some things do not need immediate decisions, but do deserve early preparation.)}

Học nền chắc, rèn thói quen tốt, giữ sức khỏe, học cách tự học và tập trách nhiệm đều là sự chuẩn bị quan trọng.
\textit{(Building strong basics, good habits, health, self-study, and responsibility are all important preparation.)}

Người chuẩn bị từ sớm thường đỡ hoảng hơn khi cơ hội hoặc thử thách tới.
\textit{(Those who prepare early are less panicked when opportunity or challenge arrives.)}

Chuẩn bị không làm em chắc chắn thắng, nhưng giúp em sẵn hơn nhiều.
\textit{(Preparation does not guarantee success, but it makes you much readier.)}
\end{truyen}

\begin{baihoc}
Chuẩn bị sớm là một cách tôn trọng tương lai của chính mình.
\textit{(Early preparation is a way of respecting your own future.)}
\end{baihoc}

\begin{ghinhoanh}
\textbf{Short English to remember:}

Preparation gives the future a stronger foundation.

\end{ghinhoanh}

\begin{tuhocnhanh}
1. Vì sao nên chuẩn bị từ sớm? \textit{Why should you prepare early?}

2. Chuẩn bị gồm những điều gì? \textit{What does preparation include?}

3. Em đang chuẩn bị điều gì cho tương lai? \textit{What are you preparing for your future?}
\end{tuhocnhanh}
\ngancach

\section{Mỗi Người Có Một Đường Riêng}
\begin{truyen}{Mỗi Người Có Một Đường Riêng}{Each Person Has a Different Path}
\chuhoa{K}{hông phải mọi người sẽ đi cùng một tốc độ hay cùng một hướng.}
\textit{(Not everyone will move at the same speed or in the same direction.)}

Có người hợp ngành này, có người hợp nghề kia, có người nở sớm, có người nở muộn.
\textit{(Some fit one field, others another; some bloom early, others later.)}

So sánh con đường mình với mọi người một cách máy móc dễ làm em mệt.
\textit{(Mechanically comparing your path with others can exhaust you.)}

Hiểu rằng mỗi người có một đường riêng giúp em bình tĩnh hơn với hành trình của mình.
\textit{(Understanding that each person has a different path helps you stay calm on your own journey.)}
\end{truyen}

\begin{baihoc}
Đi khác người khác không tự động có nghĩa là đi sai.
\textit{(Going differently from others does not automatically mean going wrong.)}
\end{baihoc}

\begin{ghinhoanh}
\textbf{Short English to remember:}

Different paths can still be right.

\end{ghinhoanh}

\begin{tuhocnhanh}
1. Vì sao mỗi người có một đường riêng? \textit{Why does each person have a different path?}

2. So sánh máy móc dễ gây điều gì? \textit{What does rigid comparison easily cause?}

3. Em cần bình tĩnh hơn với hành trình nào của mình? \textit{Where do you need more calm about your own path?}
\end{tuhocnhanh}
\ngancach

\section{Dám Mơ Nhưng Cũng Dám Làm}
\begin{truyen}{Dám Mơ Nhưng Cũng Dám Làm}{Dare to Dream and Dare to Act}
\chuhoa{Ư}{ớc mơ đẹp, nhưng ước mơ chỉ ở trong đầu thì chưa đủ.}
\textit{(Dreams are beautiful, but dreams that stay only in the mind are not enough.)}

Điều cần hơn là dám làm những phần nhỏ nhưng thật để kéo ước mơ xuống gần đời sống.
\textit{(What matters more is doing the small real things that pull a dream closer to life.)}

Dám mơ cho em hướng.
\textit{(Daring to dream gives direction.)}

Dám làm cho em con đường.
\textit{(Daring to act gives a path.)}
\end{truyen}

\begin{baihoc}
Ước mơ đẹp nhất khi nó đi cùng những việc làm cụ thể mỗi ngày.
\textit{(Dreams are most beautiful when they travel with concrete daily actions.)}
\end{baihoc}

\begin{ghinhoanh}
\textbf{Short English to remember:}

Dreams need action to become real.

\end{ghinhoanh}

\begin{tuhocnhanh}
1. Vì sao dám mơ thôi chưa đủ? \textit{Why is dreaming alone not enough?}

2. Điều gì kéo ước mơ gần lại? \textit{What brings a dream closer?}

3. Em có việc làm nào đang đi cùng ước mơ của mình? \textit{What action is walking with your dream?}
\end{tuhocnhanh}
\ngancach

\section{Tương Lai Có Thể Thay Đổi}
\begin{truyen}{Tương Lai Có Thể Thay Đổi}{The Future Can Change}
\chuhoa{T}{ương lai không phải lúc nào cũng đi đúng theo bản nháp đầu tiên em từng nghĩ.}
\textit{(The future does not always follow the first draft you once imagined.)}

Có những hướng sẽ đổi, có những cánh cửa mở ra muộn, có những điều em chỉ hiểu sau khi đi thêm một đoạn.
\textit{(Some directions change, some doors open late, and some things are understood only after more distance.)}

Biết rằng tương lai có thể thay đổi giúp em mềm hơn với đời sống.
\textit{(Knowing the future can change helps you stay more flexible in life.)}

Điều đó không xấu; nhiều khi đó là cách mình tìm được chỗ phù hợp hơn.
\textit{(That is not bad; often it is how we find a better-fitting place.)}
\end{truyen}

\begin{baihoc}
Tương lai thay đổi không luôn là mất hướng; đôi khi đó là quá trình tìm ra hướng đúng hơn.
\textit{(A changing future is not always losing direction; sometimes it is finding a better one.)}
\end{baihoc}

\begin{ghinhoanh}
\textbf{Short English to remember:}

Changing direction can still be wise.

\end{ghinhoanh}

\begin{tuhocnhanh}
1. Vì sao tương lai có thể thay đổi? \textit{Why can the future change?}

2. Điều đó có luôn xấu không? \textit{Is that always bad?}

3. Em có đang sợ việc đổi hướng quá mức không? \textit{Are you too afraid of changing direction?}
\end{tuhocnhanh}
\ngancach

\section{Giữ Hy Vọng}
\begin{truyen}{Giữ Hy Vọng}{Keep Hope}
\chuhoa{H}{y vọng không làm mọi khó khăn biến mất, nhưng nó giúp em còn lý do để tiếp tục.}
\textit{(Hope does not erase hardship, but it gives you a reason to continue.)}

Một học sinh mất hy vọng thường rất khó giữ nỗ lực lâu dài.
\textit{(A student who loses hope struggles to keep effort alive for long.)}

Giữ hy vọng không phải là tự lừa mình rằng mọi thứ sẽ dễ.
\textit{(Keeping hope is not pretending everything will be easy.)}

Nó là tin rằng mình vẫn còn một con đường để đi.
\textit{(It is believing there is still a path to walk.)}
\end{truyen}

\begin{baihoc}
Hy vọng là một ngọn đèn nhỏ, nhưng nhiều khi đủ để em không bỏ đường giữa tối.
\textit{(Hope is a small lamp, but often enough to keep you from losing the road in darkness.)}
\end{baihoc}

\begin{ghinhoanh}
\textbf{Short English to remember:}

Hope keeps the road visible.

\end{ghinhoanh}

\begin{tuhocnhanh}
1. Vì sao cần giữ hy vọng? \textit{Why do you need hope?}

2. Hy vọng có phải là tự lừa mình không? \textit{Is hope just self-deception?}

3. Điều gì đang giúp em còn hy vọng? \textit{What is helping you keep hope?}
\end{tuhocnhanh}
\ngancach

\section{Đi Từng Năm Một}
\begin{truyen}{Đi Từng Năm Một}{Take the Future Year by Year}
\chuhoa{T}{ương lai rộng quá đôi khi làm học sinh ngợp.}
\textit{(The future can feel so wide that it overwhelms students.)}

Lúc ấy, thay vì ôm hết nhiều năm vào một lúc, em có thể học cách đi từng năm, từng chặng, từng mục tiêu gần trước.
\textit{(At such times, instead of carrying many years at once, you can move year by year, stage by stage, and focus first on nearer goals.)}

Cách nghĩ này làm tương lai bớt nặng hơn.
\textit{(This way of thinking makes the future less heavy.)}

Nó cũng giúp em sống thực tế hơn với hiện tại.
\textit{(It also helps you live more realistically in the present.)}
\end{truyen}

\begin{baihoc}
Không cần gánh cả tương lai trong một ngày; đi từng chặng vẫn là đi tới.
\textit{(You do not need to carry the whole future in one day; going stage by stage is still moving forward.)}
\end{baihoc}

\begin{ghinhoanh}
\textbf{Short English to remember:}

Take the future one stage at a time.

\end{ghinhoanh}

\begin{tuhocnhanh}
1. Vì sao nên đi từng năm một? \textit{Why should you take the future year by year?}

2. Cách nghĩ này giúp gì? \textit{How does this mindset help?}

3. Chặng gần nhất em cần tập trung là gì? \textit{What is the nearest stage you need to focus on?}
\end{tuhocnhanh}
\ngancach

\section{Tương Lai Không Cần Hoàn Hảo}
\begin{truyen}{Tương Lai Không Cần Hoàn Hảo}{The Future Does Not Need to Be Perfect}
\chuhoa{N}{hiều học sinh lo vì tưởng tương lai của mình phải thật đẹp, thật đúng, thật trơn tru mới là ổn.}
\textit{(Many students worry because they think their future must be perfect, smooth, and exactly right.)}

Nhưng một tương lai có những đoạn quanh co, những lần sửa hướng và những ngày chưa như ý vẫn có thể là một tương lai rất đáng sống.
\textit{(But a future with turns, corrections, and imperfect days can still be deeply worth living.)}

Điều quan trọng không phải là không có vết gập ghềnh.
\textit{(What matters is not the absence of rough spots.)}

Điều quan trọng là em vẫn giữ được hướng sống tử tế và tiếp tục đi.
\textit{(What matters is that you keep a kind direction and continue moving.)}
\end{truyen}

\begin{baihoc}
Tương lai đủ tốt không cần phải hoàn hảo; nó chỉ cần là một con đường em đi bằng sự chân thành và nỗ lực.
\textit{(A good enough future does not need perfection; it needs a path you walk with sincerity and effort.)}
\end{baihoc}

\begin{ghinhoanh}
\textbf{Short English to remember:}

An imperfect future can still be a good one.

\end{ghinhoanh}

\begin{tuhocnhanh}
1. Vì sao tương lai không cần hoàn hảo? \textit{Why does the future not need to be perfect?}

2. Điều quan trọng hơn sự trơn tru là gì? \textit{What matters more than smoothness?}

3. Em có đang đòi tương lai của mình quá hoàn hảo không? \textit{Are you demanding too much perfection from your future?}
\end{tuhocnhanh}
\ngancach